\documentclass{article}

% Prefer a small page width for the demo
\usepackage[paperwidth=15cm]{geometry}

% Use babel with Arabic, Hindi and English languages 
% English is loaded last, making it the default language
\usepackage[hindi,english]{babel}
\babelprovide[import=ar]{arabic}
% Set Roman font for Arabic
\babelfont[arabic]{rm}{Scheherazade}

% Set Roman and sans serif fonts for English
\babelfont{rm}{Noto Serif}
\babelfont{sf}{Noto Sans}

% Set up Roman and san serif fonts for Hindi
\babelfont[hindi]{rm}[Language=Default]{Noto Serif Devanagari}
\babelfont[hindi]{sf}[Language=Default]{Noto Sans Devanagari}

% Because we want to use the "dvipsnames" option to 
% access additional named colors, we must load xcolor
% before luacolor (see the luacolor package documentation)
\usepackage[dvipsnames]{xcolor}
\usepackage{luacolor}

% Create a convenience command to typeset Hindi
\newcommand\hinditext[1]{\foreignlanguage{hindi}{#1}}

% Create a convenience command to typeset Arabic
\newcommand\arabictext[1]{\foreignlanguage{arabic}{#1}}

\usepackage{hyperref}
\hypersetup{
	colorlinks=true,
	urlcolor=cyan
}
\begin{document}
	% The Noto fonts benefit from a larger line spacing
	\setlength{\baselineskip}{14bp}
	
	\section{Colorizing Hindi text}
	Google translates the Hindi word \textsf{\hinditext{किंकर्तव्यविमूढ़}} as
	``bewildered''. By using the \texttt{luacolor} package it's possible to colorize glyphs within Hindi text; for example \hinditext{किंक\textcolor{purple}{र्तव्यव}\textcolor{green}{िमूढ़}}. You can also add color to the whole word: \hinditext{\textcolor{blue}{किंकर्तव्यविमूढ़}}.
	
	\section{Colorizing Arabic text}
	% Create a custom environment to
	% typeset colored Arabic text
	\newenvironment{colorarabic}
	{% Typeset Arabic using a larger font 
		\fontsize{30}{30}
		% Set right-to-left paragraph and text directions
		\pardir TRT\textdir TRT}
	{}
	
	Here is some colorized Arabic text:  \begin{colorarabic}
		\textcolor{red}{\arabictext{هَذَا}} \arabictext{نَصٌّ عَرَب}\textcolor{green}{\arabictext{ِ}}\arabictext{ي}\textcolor{blue}{\arabictext{ٌّ}}
	\end{colorarabic}
	\section{Colorizing diacritics}
	This example is based on \href{https://tex.stackexchange.com/questions/698933/color-breaks-diacritics-stacking}{code from tex.stackexchange}.
	
	\vspace{12pt}
	\bgroup
	\newcommand{\emptydiacritic}{\char"034F}
	\fontsize{60}{60}\selectfont
	á̀̐{\color{blue}a\emptydiacritic\color{green}́\color{red}̀\color{magenta}̐
		\egroup
	\end{document}